\section{What even an ideal agent cannot know}\label{sec:ideal}

An agent with unlimited computing power escapes \cref{thm:ud-gold}. In a deterministic
computable world every divergent rival is refuted after finitely many observations
(\cref{ud:rem-transient}), and enumeration with an oracle for totality converges on the
true behavior (\cref{sec:discover}). Two things remain even for such an agent: rival
\emph{formulations} of one final theory, and rival \emph{chance hypotheses}, which leave
the final theory itself undetermined.

\begin{theorem}[T5: Permanent underdetermination]\label{thm:t5-permanent}
\leavevmode
\begin{enumerate}[label=(\alph*)]
  \item \emph{Formulations.} If two models are extensionally equal
        (\cref{def:model})---deterministic models $e,e'$ with $\ph{e}=\ph{e'}$, or
        stochastic models $i,j$ with $\nu_i=\nu_j$---then they assign the same
        outcome, or the same probability, to every possible record. No method of any kind, computable or not, separates
        them.
  \item \emph{Chance hypotheses.} Let $\mu$ and $\mu'$ be distinct computable
        measures on $\Seq$ that are mutually absolutely continuous, that is, agree on
        which events have probability zero: $\mu(B)=0$ if and only if $\mu'(B)=0$, for
        every Borel set $B$. Let $L$ be any method, computable or not, mapping finite outcome
        strings to models. If $L$ converges to a model of $\mu$ with $\mu$-probability
        $1$, then it converges to a model of $\mu'$ with $\mu'$-probability $0$. In
        particular, no method identifies the correct one of $\mu,\mu'$ with probability
        $1$ under both.
\end{enumerate}
\end{theorem}

\begin{proof}
(a) A record's outcomes are fixed by its settings and the model's behavior on them
(\cref{def:model}); extensionally equal models agree on both, so every method receives
the same input whichever of them is true.

(b) For a set $I$ of models (a countable set of indices), the event that $L$ converges
to a member of $I$ is
\[
  A_I=\bigcup_{N}\,\bigcup_{i\in I}\,\bigcap_{n\ge N}\{\omega\in\Seq : L(\omega_{\le n})=i\}.
\]
Each set $\{\omega : L(\omega_{\le n})=i\}$ depends only on the first $n$ symbols, so it
is a finite union of cylinders; the unions and intersections are countable. Hence $A_I$
is Borel whatever $L$ is. Let $I$ and $I'$ be the sets of models of
$\mu$ and of $\mu'$, that is, of indices $i$ with $\nu_i=\mu$ and with $\nu_i=\mu'$. They are disjoint because $\mu\neq\mu'$, so $A_I\cap A_{I'}=\emptyset$.
If $\mu(A_I)=1$ then $\mu(A_{I'})=0$, and by absolute continuity $\mu'(A_{I'})=0$.
\end{proof}

The argument uses no property of $\mu,\mu'$ beyond mutual absolute continuity. The
restriction to computable measures loses nothing: a probability measure that has a
stochastic model is lower semicomputable, and a lower semicomputable probability
measure is computable, since $\nu(r)=1-\sum_{r'\ne r,\,|r'|=|r|}\nu(r')$ is also upper
semicomputable \citep{li-vitanyi2019}.

A method that uses internal randomness is covered as well, provided each conjecture
depends on $\omega_{\le n}$ and on finitely many of the method's random bits $\rho$
(possibly a number of them that depends on the bits read). Then the success event $B$
is a Borel subset of $\Seq\times\{0,1\}^{\infty}$, and (b) applies to the products
$\mu\times\lambda$ and $\mu'\times\lambda$, where $\lambda$ is the uniform measure on the
random bits. These products are again mutually absolutely continuous: by Fubini's
theorem, $(\mu\times\lambda)(B)=\int\mu(B^{\rho})\,d\lambda(\rho)$, which vanishes if and
only if $\mu(B^{\rho})=0$ for $\lambda$-almost every $\rho$, hence if and only if
$\mu'(B^{\rho})=0$ for $\lambda$-almost every $\rho$, hence if and only if
$(\mu'\times\lambda)(B)=0$.

\paragraph{Instances of (a).}
The trivial instance, as an illustration only, is the padding lemma: every behavior has
infinitely many programs (\cref{thm:ud-infinite}(b)). The serious instances are rival \emph{formulations} that
physicists take seriously.
\begin{itemize}
  \item \textbf{Dualities.} Pairs of theories with very different ingredients are held to
        describe the same physics. In T-duality, string theory on a circle of radius
        $R$ is equivalent to string theory on a circle of radius $\alpha'/R$
        \citep{polchinski1998}. In the AdS/CFT correspondence, string theory on an
        anti-de Sitter space is conjectured to be equivalent to a gauge theory without
        gravity on its boundary \citep{maldacena1998}. Where such a duality holds, no
        record says which side is fundamental, or whether the space has one size or the
        other.
  \item \textbf{Interpretations of quantum mechanics.} The following all assign Born-rule
        probabilities to every record:
        \begin{itemize}
          \item Copenhagen-style quantum mechanics, in which the probabilities are
                fundamental;
          \item de Broglie--Bohm theory in quantum equilibrium, which is deterministic,
                with randomness confined to the initial positions
                \citep{bohm1952,dgz1992};
          \item Everettian quantum mechanics, which is deterministic unitary evolution,
                with Born weights for branches \citep{everett1957,wallace2012}.
        \end{itemize}
        As stochastic models they are therefore extensionally equal. The question ``is
        the world fundamentally random?'' is not settled by any record.
\end{itemize}
Two candidates do not belong on the list. Objective-collapse theories \citep{grw1986}
predict small departures from standard unitary quantum predictions (suppressed
macroscopic interference, spontaneous heating), so they are rivals that records can in
principle refute, not instances of (a). Superdeterministic proposals are programmatic:
their proponents describe the models as ``still in their infancy''
\citep[Sec.~5]{hossenfelder-palmer2020}.

\paragraph{Instances of (b).}
The simplest example is two coin models. Under $\mu$ every toss lands heads with
probability $1/2$. Under $\mu'$ the first toss lands heads with probability $1/3$ and
every later toss with probability $1/2$. These are distinct chance hypotheses, since
they disagree about a definite event. Because they differ in only finitely many
conditional probabilities, each of which is strictly between $0$ and $1$, they are
mutually absolutely continuous (the density of $\mu'$ with respect to $\mu$ is $2/3$ or
$4/3$). More generally, two theories of chance are mutually absolutely continuous if
their conditional probabilities differ at only finitely many histories in total, at
each of which both conditionals lie in $(0,1)$: the density is then a finite product of
ratios bounded away from $0$ and $\infty$. For independent trials with pairwise equivalent
factors $\mu_n,\mu'_n$, Kakutani's criterion gives mutual absolute continuity exactly
when $\prod_n\int\sqrt{d\mu_n\,d\mu'_n}>0$; for Bernoulli factors with probabilities
$p_n,q_n$ bounded away from $0$ and $1$, this holds if and only if
$\sum_n(p_n-q_n)^2<\infty$ \citep{kakutani1948}. A method may still guess well: if the
first toss has heads probability $0.99$ under one hypothesis and $0.01$ under the other,
guessing from that toss is right with probability $0.99$ either way. What no method
achieves is convergence on the correct hypothesis with probability $1$ under both.
Positive results exist for restricted classes. \Citet{vitanyi-chater2017} identify a
computable measure in the limit when the data are independent draws from a
computably enumerable class of measures; for dependent data their learner returns a
measure for which the sequence is typical, and there may be several, which is the gap
(b) exhibits. Without such restrictions, algorithmic identification of probabilities
fails \citep{bienvenu2018}.

\paragraph{Scope.}
The theorem takes a physical formulation to be a model in the sense of \cref{def:model},
so that whether a formulation is fundamental is a property of the model rather than of
its predictions. Part (a) then shows that the \emph{fundamental formulation} of a theory
of everything is never settled by data: two candidates with the same predictions may
disagree about what exists (strings in one space or a gauge theory on its boundary;
particles with trajectories or a branching wave function). Part (b) reaches the final
theory itself: for theories of chance, a rival that agrees on which events have
probability zero puts the final theory, in the sense of \cref{def:final}, beyond
identification with probability 1 under both.
